\subsection{Methods, Algorithms, and Computational Artifacts}

Procedural blocks must distinguish conceptual guidance, executable steps, and claims about correctness or performance.

\begin{convention} \label{conv:procedural-environment-roles} \hlt{Procedural Environment Roles}\\
Use the following roles throughout these notes.
\begin{enumerate}[label=(\roman*)]
\setlength{\itemsep}{0pt}
\item \hlt{Proof method}: a reusable proof route, reduction, recognition criterion, or obligation checklist.
\item \hlt{Computation method}: a reusable construction, calculation, approximation, diagnostic, or test.
\item \hlt{Method}: a short convention, classification workflow, or inseparable hybrid with no clearer subtype.
\item \hlt{Process}: an organizational workflow for research, writing, review, or documentation.
\item \hlt{Algorithm}: executable pseudocode with declared inputs, outputs, control flow, and termination conditions.
\end{enumerate}
A worked instance belongs in an example, an actual proof in a proof, and implementation-specific behavior in code or an artifact record.
\end{convention}

\begin{process} \label{proc:algorithm-specification-audit} \hlt{Algorithm Specification and Audit}\\
Before publishing pseudocode, record its contract and verification obligations.
\begin{enumerate}[label=(\roman*)]
\setlength{\itemsep}{0pt}
\item \hlt{Problem}: state the task and the abstraction level at which the algorithm operates.
\item \hlt{Input}: specify types, dimensions, admissibility, representation, and malformed-input behavior.
\item \hlt{Output}: specify the returned object and a postcondition that characterizes success.
\item \hlt{State}: define initialization, mutable state, randomness, and any invariant needed for correctness.
\item \hlt{Steps}: write one unambiguous action per step with readable branches and loops.
\item \hlt{Termination}: state the stopping rule and justify termination when it is claimed.
\item \hlt{Correctness}: separate the proof of the postcondition from the pseudocode itself.
\item \hlt{Resources}: state complexity only with an input-size parameter, cost model, resource, and case type.
\item \hlt{Limits}: record numerical, memory, precision, platform, approximation, and applicability boundaries.
\end{enumerate}
Labeled steps and explanatory invariants clarify the distinction between an algorithm and program text \citep{knuth1988mathematicalwriting}.
\end{process}

\begin{principle} \label{princ:formal-informal-procedural-pair} \hlt{Formal and Informal Procedural Views}\\
Complex procedures often need complementary views. The formal view fixes inputs, outputs, updates, and stopping conditions; the informal view explains the invariant, decomposition, or design choice. Neither should duplicate the other, and only a separate proof can establish mathematical correctness.
\end{principle}

\begin{process} \label{proc:computational-artifact-record} \hlt{Computational Artifact Record}\\
For code supporting a reported result, archive the source revision, dependencies, runtime, platform, configuration, random seeds, data identifiers, licenses, commands, and expected outputs. Map each reported figure, table, or metric to the command that creates it. State unavailable components and access restrictions explicitly. Current ML reporting guidance requests exact execution instructions and enough experimental detail to interpret the main results \citep{neurips2026checklist}.
\end{process}

\begin{warning} \label{warn:algorithm-evidence-boundaries} \hlt{Algorithm Evidence Boundaries}\\
Passing tests does not establish correctness outside the tested cases, and an asymptotic bound does not predict wall-clock performance without a cost model and operating conditions. Keep specification, proof, implementation, benchmark, and interpretation as distinct evidence types.
\end{warning}
